\chapter{Vectors in \texorpdfstring{$\mathbb{R}^n$}{R\string^n}}
\label{mf:ch02-vectors-rn}

The next part of the course begins the linear algebra pillar. The basic object is a vector. Vectors will later be added, scaled, multiplied through dot products, collected into matrices, and used to describe regression geometry. Before doing any of that, we need a clear answer to a simpler question: what is a vector?

In this chapter, a vector is an ordered collection of real numbers. The entries might represent coordinates in physical space, but they might also represent inventory amounts, probabilities, measurements over time, or values of a variable in a data set. The important point is that the entries are organized and interpreted together.

\section{The Space \texorpdfstring{$\mathbb{R}^n$}{R\string^n}}
\label{mf:vectors-rn}

Let \(\mathbb{R}\) denote the set of real numbers. For a positive integer \(n\), the set \(\mathbb{R}^n\) is the set of all ordered lists of \(n\) real numbers. An element of \(\mathbb{R}^n\) is called an \emph{\(n\)-vector}. Unless stated otherwise, vectors are written as columns:
\[
\mathbf{x}
=
\begin{bmatrix}
 x_1 \\
 x_2 \\
 \vdots \\
 x_n
\end{bmatrix}
\in \mathbb{R}^n.
\]
The same vector is sometimes written on one line as
\[
\mathbf{x} = (x_1,x_2,\ldots,x_n)^T,
\]
where the superscript \(T\) indicates transpose. The transpose notation reminds us that the displayed list \((x_1,x_2,\ldots,x_n)\) is being treated as a column vector.

The number \(n\) is the \emph{dimension} or \emph{length} of the vector. A scalar is a single number; a vector is a structured collection of numbers. For example,
\[
\begin{bmatrix}
5 \\
2 \\
9
\end{bmatrix}
\in \mathbb{R}^3
\]
is a 3-vector, while \(5\in\mathbb{R}\) is a scalar.

\section{Entries, Coordinates, and Subvectors}
\label{mf:vector-entries-subvectors}

For a vector \(\mathbf{x}\in\mathbb{R}^n\), the scalar values
\[
 x_1,x_2,\ldots,x_n
\]
are called the \emph{entries}, \emph{components}, \emph{elements}, or \emph{coordinates} of \(\mathbf{x}\). The notation \(x_i\) means the \(i\)th entry of \(\mathbf{x}\).

It is often useful to build a larger vector by stacking smaller vectors. Suppose
\[
\mathbf{b}\in\mathbb{R}^m,
\qquad
\mathbf{c}\in\mathbb{R}^n,
\qquad
\mathbf{d}\in\mathbb{R}^p.
\]
Then the stacked vector
\[
\mathbf{a}
=
\begin{bmatrix}
\mathbf{b} \\
\mathbf{c} \\
\mathbf{d}
\end{bmatrix}
\]
belongs to \(\mathbb{R}^{m+n+p}\). In this situation, \(\mathbf{b}\), \(\mathbf{c}\), and \(\mathbf{d}\) are \emph{subvectors} or \emph{slices} of \(\mathbf{a}\).

Colon notation is often used to describe consecutive entries. If \(r\leq s\), then
\[
\mathbf{a}_{r:s} = (a_r,a_{r+1},\ldots,a_s)^T.
\]
For the stacked vector above,
\[
\mathbf{b}=\mathbf{a}_{1:m},
\qquad
\mathbf{c}=\mathbf{a}_{m+1:m+n},
\qquad
\mathbf{d}=\mathbf{a}_{m+n+1:m+n+p}.
\]
This notation will later be useful for selecting portions of vectors, rows of data, and groups of variables.

\section{Indexing Conventions}
\label{mf:vector-indexing}

In the mathematical notation used in this course, vectors are usually \emph{1-indexed}. That means the first entry of \(\mathbf{x}\) is \(x_1\), the second entry is \(x_2\), and the final entry is \(x_n\).

Some programming languages, including Python and C/C++, are \emph{0-indexed}. In those languages, the first entry of a stored vector or array has index 0. This difference is not a mathematical disagreement; it is a convention. When translating between mathematics and code, always check whether an index begins at 0 or 1.

\section{Special Vectors}
\label{mf:special-vectors}

Several vectors occur so often that they receive special notation.

\paragraph{The zero vector.}
\label{mf:zero-vector}
The \emph{zero vector} in \(\mathbb{R}^n\), written \(\mathbf{0}_n\) or simply \(\mathbf{0}\) when the dimension is clear, contains only zeros:
\[
\mathbf{0}_n=
\begin{bmatrix}
0\\0\\\vdots\\0
\end{bmatrix}.
\]

\paragraph{The ones vector.}
\label{mf:ones-vector}
The \emph{ones vector} in \(\mathbb{R}^n\), written \(\mathbf{1}_n\) or simply \(\mathbf{1}\), contains only ones:
\[
\mathbf{1}_n=
\begin{bmatrix}
1\\1\\\vdots\\1
\end{bmatrix}.
\]
The ones vector is especially important later because it represents an intercept column in regression and helps express sample sums and sample means.

\paragraph{Standard basis vectors.}
\label{mf:standard-basis-vectors}
For \(i=1,\ldots,n\), the \(i\)th \emph{standard basis vector} \(\mathbf{e}_i\in\mathbb{R}^n\) has a 1 in position \(i\) and 0 in every other position. For example, in \(\mathbb{R}^3\),
\[
\mathbf{e}_1=\begin{bmatrix}1\\0\\0\end{bmatrix},
\qquad
\mathbf{e}_2=\begin{bmatrix}0\\1\\0\end{bmatrix},
\qquad
\mathbf{e}_3=\begin{bmatrix}0\\0\\1\end{bmatrix}.
\]
These vectors will become important when we discuss linear combinations and bases.

\section{Sparse and Dense Vectors}
\label{mf:sparse-vectors}

A vector is called \emph{sparse} if many of its entries are zero. A vector that is not sparse is often called \emph{dense}. The \emph{sparsity pattern} of a vector is the set of positions where the entries are nonzero.

For example,
\[
\begin{bmatrix}
0\\0\\7\\0\\0\\2
\end{bmatrix}
\]
is sparse because only two of its six entries are nonzero. We sometimes write \(\operatorname{nnz}(\mathbf{x})\) for the number of nonzero entries of \(\mathbf{x}\).

Sparsity is not merely a technical detail. If a vector records the quantity of many different resources in a warehouse, sparsity means that most resource types are absent. If a vector records words in a document, sparsity means that most words in the vocabulary do not appear in that document. Recognizing sparsity can simplify computation and interpretation.

\section{Vectors as Data and Measurements}
\label{mf:vectors-as-data}

A vector should not be thought of only as a point in physical space. That picture is useful in some settings, but the broader idea is more important:

\begin{quote}
A vector is an ordered collection of quantities that pertain to the same object, system, or concept.
\end{quote}

Here are several common interpretations.

\begin{center}
\begin{tabular}{p{0.27\textwidth} p{0.58\textwidth}}
\hline
Vector interpretation & Example \\
\hline
Location & \((\text{longitude},\text{latitude},\text{elevation})\in\mathbb{R}^3\) \\
Displacement & \((\Delta\text{longitude},\Delta\text{latitude},\Delta\text{elevation})\in\mathbb{R}^3\) \\
Color & RGB values such as \((1,0,0)\) for red or \((0,1,0)\) for green \\
Quantities & Amounts of different resources, products, or inventory classes \\
Time series & Values of a signal or measurement at successive times \\
Data record & Measurements or features describing one observation \\
\hline
\end{tabular}
\end{center}

The same mathematical object can therefore support different interpretations. In operations research and data analysis, this flexibility is essential. A vector might describe a ship's position, a unit's inventory, a time series of sensor readings, or one row of a data table.

\paragraph{Data Analysis bridge.}
In the Data Analysis textbook, a feature vector is one observation's ordered collection of input measurements. A response vector is the collection of observed outputs over many observations. Later, a residual vector will collect the leftover errors from a fitted model. All of these are instances of the same basic object: a vector in some \(\mathbb{R}^n\).

\section{Probability Vectors and Time Series}
\label{mf:probability-vectors}
\label{mf:time-series-vectors}

Two examples deserve special attention because they recur often in modeling.

A \emph{probability vector} records the probabilities of a finite set of outcomes. For example, for a fair six-sided die,
\[
\mathbf{w}=\left(\frac16,\frac16,\frac16,\frac16,\frac16,\frac16\right)^T.
\]
For \(\mathbf{w}\) to be a valid probability vector, its entries must be nonnegative and must sum to 1:
\[
 w_i\geq 0 \quad \text{for all } i,
 \qquad
 \sum_{i=1}^n w_i=1.
\]
Probability vectors will later connect naturally to weighted averages and expectations.

A \emph{time-series vector} records the value of some quantity at ordered time points. For example, an audio signal can be represented as a vector whose entries record acoustic pressure at equally spaced times. When a vector represents a time series, it is natural to plot the entry \(x_i\) against the time index \(i\).

\section{Chapter Summary}
\label{mf:vectors-summary}

Carry forward these vector ideas:
\begin{itemize}
    \item \(\mathbb{R}^n\) is the set of all length-\(n\) real vectors; vectors have entries, components, or coordinates.
    \item Subvectors such as \(\mathbf{a}_{r:s}\) select consecutive entries, and mathematical indexing usually begins at 1.
    \item The zero vector, ones vector, and standard basis vectors are recurring building blocks.
    \item Sparse vectors have many zero entries; dense vectors do not.
    \item In this course, vectors are not limited to physical arrows. They can represent records, quantities, probabilities, time series, or data columns.
\end{itemize}

The next chapter develops vector arithmetic and linear combinations.
